\section{矩估计让样本特征匹配总体特征}
\label{sec:11-Mathematical-Statistics/02-Point-Estimation/02}

排队仿真需要一个到达间隔的分布。日志里有四十万条时间戳，差分之后直方图右偏、长尾，Gamma 分布是个合理的候选。问题是形状参数和尺度参数取多少。

打开文档，最大似然那条路要解一个含 digamma 函数的方程，得配牛顿迭代，还得处理初值和收敛判据。另一条路只要两行：算样本均值 \(\bar X=4.2\) 秒，算样本方差 \(S^2=31.8\)，然后

\[
\hat\vartheta=\frac{S^2}{\bar X}=7.57,\qquad \hat k=\frac{\bar X^2}{S^2}=0.555 .
\]

没有迭代，没有初值，没有收敛判据。仿真跑起来了。这一节讲这两行公式从哪儿来，以及它省掉的那些东西是白省的还是要还的。

\subsection{矩方程是一个联立方程组}

参数 \(\theta\in\R^k\) 有 \(k\) 个分量。总体的第 \(j\) 阶原点矩是参数的函数，\(m_j(\theta)=\E_\theta[X^j]\)；样本的第 \(j\) 阶原点矩直接从数据算，\(\hat m_j=\frac1n\sum_i X_i^j\)。矩估计把这两列量对齐：

\[
m_j(\hat\theta)=\hat m_j,\qquad j=1,\dots,k .
\]

\(k\) 个方程，\(k\) 个未知数，通常恰好定解。整个方法就这一句话，剩下的是代数。

Gamma 分布用形状 \(k\) 与尺度 \(\vartheta\) 参数化时，\(\E[X]=k\vartheta\)、\(\Var(X)=k\vartheta^2\)。前两阶矩换成均值和方差写起来更顺手，因为方差是前两阶原点矩的可逆变换，匹配 \(\set{m_1,m_2}\) 与匹配 \(\set{\text{均值},\text{方差}}\) 得到的是同一个解。于是

\[
\begin{aligned}
\hat k\hat\vartheta&=\bar X,\\
\hat k\hat\vartheta^2&=S^2,
\end{aligned}
\qquad\Longrightarrow\qquad
\hat\vartheta=\frac{S^2}{\bar X},\quad \hat k=\frac{\bar X^2}{S^2}.
\]

骨架一步：两式相除消掉 \(\hat k\)，得 \(\hat\vartheta\)；回代得 \(\hat k\)。条件对账：需要 \(\bar X>0\) 与 \(S^2>0\)，前者由 Gamma 的支持集保证（只要样本非空），后者在所有观测值相同时失效。这里的 \(S^2\) 用 \(n\) 作分母还是 \(n-1\)，矩方程本身没有偏好——匹配的是二阶总体矩，除以 \(n\) 是它的原生形式。

Bernoulli 只有一个参数，\(\E[X]=\theta\)，一阶矩方程直接给出 \(\hat\theta=\bar X\)。这个模型上矩估计和最大似然给出同一个答案，所以它看不出两条路线的区别。区别要到别的模型上才显形。

\subsection{很多份不同的数据压出同一对矩}

矩估计的动作，是把整份数据压成 \(k\) 个数，再从这 \(k\) 个数里反解参数。压缩是它省事的来源，也是它全部毛病的来源。

\begin{center}
\begin{tikzpicture}[font=\small,>=stealth]
  % ---- 左：两份形状不同的数据，同一对矩 ----
  \node[font=\footnotesize,black!70] at (1.50,2.05) {两份数据};
  \foreach \i/\h in {0/0.10, 1/0.28, 2/0.58, 3/0.92, 4/1.15, 5/1.15, 6/0.92, 7/0.58, 8/0.28, 9/0.10}
    {\fill[black!20] ({0.30+0.24*\i},0.55) rectangle ({0.54+0.24*\i},{0.55+0.95*\h});
     \draw[black!45] ({0.30+0.24*\i},0.55) rectangle ({0.54+0.24*\i},{0.55+0.95*\h});}
  \draw[black!55] (0.24,0.55) -- (2.76,0.55);
  \foreach \i/\h in {0/0.08, 1/0.38, 2/0.80, 3/0.55, 4/0.22, 5/0.22, 6/0.55, 7/0.80, 8/0.38, 9/0.08}
    {\fill[black!20] ({0.30+0.24*\i},-1.95) rectangle ({0.54+0.24*\i},{-1.95+0.95*\h});
     \draw[black!45] ({0.30+0.24*\i},-1.95) rectangle ({0.54+0.24*\i},{-1.95+0.95*\h});}
  \draw[black!55] (0.24,-1.95) -- (2.76,-1.95);
  \draw[dashed,black!60] (1.50,-2.30) -- (1.50,1.80);
  \node[font=\footnotesize,black!60,anchor=north] at (1.50,-2.32) {\(\bar X\) 相同};
  \draw[<->,black!60] (0.88,0.28) -- (2.12,0.28);
  \draw[<->,black!60] (0.88,-2.22) -- (2.12,-2.22);
  \node[font=\footnotesize,black!60,anchor=west] at (2.16,0.28) {\(S^2\) 也相同};
  % ---- 箭头 ----
  \draw[->,black!65] (3.10,-0.70) -- (3.95,-0.70);
  \node[font=\footnotesize,black!60,anchor=south] at (3.52,-0.62) {矩方程};
  % ---- 右：矩平面与可行域 ----
  \begin{scope}[shift={(4.30,-2.05)}]
    \fill[black!10] (0,0) -- (2.45,2.45) -- (0,2.45) -- cycle;
    \draw[->,black!65] (0,0) -- (2.90,0) node[font=\footnotesize,black!70,right] {\(\bar X\)};
    \draw[->,black!65] (0,0) -- (0,2.90) node[font=\footnotesize,black!70,above] {\(S^2\)};
    \draw[black!55] (0,0) -- (2.45,2.45);
    \node[font=\footnotesize,black!55,anchor=north,rotate=45] at (2.10,2.10) {\(S^2=\bar X\)};
    \fill[black!80] (0.85,1.95) circle (1.8pt);
    \node[font=\footnotesize,black!60,anchor=north] at (0.85,1.80) {\(\hat r>0\)};
    \draw[black!85,line width=0.9pt] (1.78,0.92) -- (2.02,1.16);
    \draw[black!85,line width=0.9pt] (2.02,0.92) -- (1.78,1.16);
    \node[font=\footnotesize,black!60,anchor=north] at (1.90,0.80) {\(\hat r<0\)};
  \end{scope}
\end{tikzpicture}

\emph{左边两份数据的前两阶矩完全一致，矩方程无从分辨；右边的矩平面上，落在分界线下方的样本会把参数解到合法区域之外。}
\end{center}

左边那两个直方图，一个单峰对称，一个明显双峰，可它们的样本均值和样本方差一模一样。矩估计只读这两个数，交出的参数因此完全相同。似然函数不会这样——它把每个观测值单独代进密度里，双峰的那份数据在单峰模型下的似然会低得多。这是效率差距的来源：被丢掉的信息本来可以用来把参数钉得更死。

右边那张图说的是另一件事。样本矩落在平面上的位置由数据决定，而模型能够生成的矩只覆盖平面的一部分。当样本点落到可行区域外面，矩方程仍然有解，只是这个解不在参数空间里。下一小节就是这种情形。

\subsection{解总是解得出来，未必落在参数空间里}

监控系统要给某个接口每分钟的错误数拟合负二项分布，用来设报警阈值。负二项的均值和方差满足 \(\mu=r(1-p)/p\)、\(\sigma^2=\mu/p\)，反解得到

\[
\hat p=\frac{\bar X}{S^2},\qquad \hat r=\frac{\bar X^2}{S^2-\bar X}.
\]

这一周的数据是 \(\bar X=3.4\)、\(S^2=2.9\)。代进去，\(\hat p=1.17\)，\(\hat r=-23.1\)。概率大于一，形状参数是负的。求解器不会报错，它只是做了两次除法；崩掉的是下游那个按参数采样的仿真模块，堆栈里看不出任何和统计有关的线索。

病因不在代码。负二项分布的方差恒大于均值，这是模型自带的结构性约束——它只能描述过离散的计数。这一周的错误数是欠离散的，数据落在了模型能够生成的矩的区域之外。矩方程对此一无所知：它只是把两个数代进两个公式，公式不知道自己有定义域。

同一个毛病还有一个经典版本。\(\Uniform(0,\theta)\) 的一阶矩是 \(\theta/2\)，矩估计给出 \(\hat\theta=2\bar X\)。抽到 \(0.3,\ 0.8,\ 1.9,\ 0.7,\ 8.2\) 这五个数，\(\bar X=2.38\)，估计值是 \(4.76\)。可样本里明明有一个 \(8.2\)。这个估计断言，一条已经发生过的观测在模型下概率为零。分布的支持集是关于参数的最强信息之一，矩估计从平均值那道窄门进去，完全绕开了它。

最大似然在这个模型上给出 \(\hat\theta=\max_i X_i\)，它由支持集约束直接确定，不可能产生和数据矛盾的参数。

\subsection{能算出来和算得好是两件事}

上面两个例子还可以辩解说是模型选错了。就算模型完全正确、参数解也完全合法，矩估计和最大似然之间仍有一道差距。

还看 \(\Uniform(0,\theta)\)，两个估计量都一致，都能算，都不需要迭代。算一下它们的均方误差。矩估计 \(2\bar X\) 无偏，方差是

\[
\Var(2\bar X)=\frac{4}{n}\cdot\frac{\theta^2}{12}=\frac{\theta^2}{3n} .
\]

最大似然 \(\max_i X_i\) 有偏，期望 \(n\theta/(n+1)\)，把偏差平方和方差加起来化简得

\[
\MSE\bigl(\max_i X_i\bigr)=\frac{2\theta^2}{(n+1)(n+2)} .
\]

一个以 \(1/n\) 的速率下降，另一个以 \(1/n^2\)。取 \(n=100\)、\(\theta=1\)：前者是 \(3.3\times 10^{-3}\)，后者是 \(1.9\times 10^{-4}\)，差 \(17\) 倍。换个说法更刺眼——矩估计要收到 \(1717\) 个样本，才能达到最大似然用 \(100\) 个样本的精度。

这个例子极端在于支持集参数的特殊性，但方向是普遍的：在正则模型里，最大似然渐近达到 Cramér--Rao 下界，而矩估计一般达不到，两者的渐近方差之比就是效率损失，见\hyperref[sec:11-Mathematical-Statistics/03-Sufficient-Statistics-and-Information/03]{``Cramér--Rao 界给无偏估计设下限''}。开头那个 Gamma 拟合，形状参数 \(\hat k=0.555\) 落在小形状区，恰好是矩估计效率最差的地方；仿真只关心分布的大致形态时这没关系，要拿这个 \(\hat k\) 去做假设检验就不合适了。

还有一个实践上更常遇到的坑：高阶矩的收敛可以慢到没有意义。四阶样本矩由少数几个极端观测主导，重尾数据上 \(n=10^4\) 都不足以让它稳定下来。名义上的一致性说``总会收敛''，它不说要等到什么时候。Cauchy 分布更干脆，一阶矩根本不存在，矩估计的第一步就踏空。

\subsection{矩条件多于参数时要加权}

有时手上的矩条件不止 \(k\) 个。多几条约束听起来是好事，可 \(k\) 个方程已经把 \(k\) 个未知数定死了，再加方程，一般没有任何参数能同时精确满足全部。

广义矩估计的办法是从``解方程''退回到``让残差尽量小''：

\[
\hat\theta=\argmin_\theta\ \hat g(\theta)^{\trans}W\hat g(\theta),
\]

其中 \(\hat g(\theta)\) 是各条矩条件的样本残差堆成的向量，\(W\succeq 0\) 是权重矩阵。权重不是装饰——不同矩条件的量纲和噪声水平差得远，等权相加等于让方差最大的那条主导结果。最优权重是残差协方差的逆，它本身要先估出来，于是有了两步法。

多出来的矩条件顺带提供了一个检验：模型正确时加权残差不该大，大了说明模型和数据在某个方向上对不上。这个检验的结论只在模型族内部有效，它能说``这个设定有问题''，说不出问题出在现实的哪一处。

\subsection{矩估计不是最大似然的粗糙版本}

把矩估计当作过渡方案会低估它。它有几处不可替代的用途：理论矩与参数的关系显式可解时，一行代码换来一个可用的答案；似然函数写不出来或代价太高时（隐式模型、模拟模型、复杂依赖结构），矩条件可能是唯一能落地的接口；只想依赖有限几条总体特征、不愿对整个分布形状表态时，矩条件天然是半参数的；迭代算法需要一个像样的初值时，矩估计几毫秒就给出来。经济学的 Euler 方程、流行病学的再生数关系，本身就是以矩条件的形式写下来的。

让出的东西这一节已经列清了：丢弃分布形状里的信息，效率一般低于最大似然；不尊重参数空间的硬约束，可能解出非法值；高阶矩被尾部绑架。选择矩估计，等于声明``我相信这几个总体特征已经装下了参数的全部信息''。这句声明有时对，有时不对，值得写进文档而不是留在代码里。

矩估计的构造原理是对齐汇总量。还有一条完全不同的原理：不去汇总，而是直接问每个参数值把手上这份数据解释得有多好，然后挑最好的那个。下一节从这个问题开始。
